\section*{Chapitre \ref{ch:b3:heterocycles} --- Chimie des hétérocycles}

\begin{solution}{exo:b3:heterocycles:1}
Oxazole~: O en apporte deux (un doublet dans le système $\pi$, un dans le plan), N un, les
trois carbones trois~: 6. Thiazole~: de même avec S~: 6. Pyrimidine~: un
par chacun des six atomes, les deux doublets des azotes étant dans le
plan~: 6. Purine~: neuf atomes, l'azote N--H en apportant deux et les
autres un chacun~: 10, un nombre de Hückel
($4n + 2$, $n = 2$).
\end{solution}

\begin{solution}{exo:b3:heterocycles:2}
Pipéridine (11.1) $>$ DMAP (9.6) $>$ imidazole (7.0) $>$ pyridine (5.2) $\gg$ pyrrole (environ $-4$,
protoné sur un carbone).
\end{solution}

\begin{solution}{exo:b3:heterocycles:3}
2-Bromopyrrole (le pyrrole est si réactif qu'il se polybrome facilement)~;
2-bromofurane~; 3-bromoindole.
\end{solution}

\begin{solution}{exo:b3:heterocycles:4}
2,5-Diméthylfurane~; 1,2,5-triméthylpyrrole~; 2,5-diméthylthiophène.
\end{solution}

\begin{solution}{exo:b3:heterocycles:5}
2-Méthoxypyridine. L'addition du méthanolate en C2 place la charge
négative du \omterm{def:b3:heterocycles:snar}{complexe de Meisenheimer} sur l'azote~; en C3, elle ne peut se
répartir que sur des carbones, si bien que l'intermédiaire, et l'état de
transition qui le précède, sont beaucoup plus hauts en énergie.
\end{solution}

\begin{solution}{exo:b3:heterocycles:6}
2-Aminopyridine (sous forme de son sel de sodium, puis protonée lors du
traitement). L'ion amidure s'additionne en C2, ce qui donne un adduit $\sigma$ anionique portant H et $\ce{NH2}$ sur C2 et la charge
négative sur l'azote du cycle~; il perd un hydrure, qui arrache un proton
au groupe amino et part sous forme de $\ce{H2}$.
\end{solution}

\begin{solution}{exo:b3:heterocycles:7}
L'oxygène du N-oxyde rend un doublet au cycle~: des formes mésomères
portent une charge négative en C2, C4 et C6. Les électrophiles attaquent C4
(C2 est proche de l'azote chargé positivement). L'oxygène est ensuite
éliminé par le trichlorure de phosphore, ce qui donne la 4-nitropyridine.
\end{solution}

\begin{solution}{exo:b3:heterocycles:8}
2-Hydroxypyridine (un OH sur C2) et 2-pyridone (N--H et C=O). La pyridone conserve six
électrons $\pi$ dans le cycle~: sa forme mésomère zwitterionique, avec $\mathrm{O^-}$ et un
$\mathrm{N^+}$ qui partage son doublet, est un pyridinium-olate~; c'est un amide dont le
doublet de l'azote complète le sextet, et les fortes liaisons hydrogène
de C=O et de N--H la favorisent dans l'eau.
\end{solution}

\begin{solution}{exo:b3:heterocycles:9}
Le 2-nitrobenzaldéhyde, deux équivalents d'acétoacétate de méthyle et
l'ammoniac~: l'inhibiteur calcique nifédipine.
\end{solution}

\begin{solution}{exo:b3:heterocycles:10}
L'ène-hydrazine orientée vers le groupe $\ce{CH2}$ (interne, plus substituée) donne le
2,3-diméthylindole~; celle orientée vers le groupe méthyle (terminale)
donne le 2-éthylindole. Les acides faibles donnent surtout le premier, par
l'ène-hydrazine la plus stable~; des acides plus forts augmentent la
proportion de 2-éthylindole.
\end{solution}

\begin{solution}{exo:b3:heterocycles:11}
En $\mathrm S_\mathrm NAr$, l'étape cinétiquement déterminante est l'addition, qui ne rompt pas la liaison C--X~;
le fluor, plus électronégatif, stabilise davantage l'état de transition
anionique et le \omterm{def:b3:heterocycles:snar}{complexe de Meisenheimer}. En $\mathrm S_\mathrm N2$, la liaison C--X se rompt dans
l'état de transition, et la liaison C--I, faible, part le plus vite.
\end{solution}

\begin{solution}{exo:b3:heterocycles:12}
Pyridine~: trois signaux dans le rapport 2\,:\,1\,:\,2, le premier vers \qty{8.6}{ppm}. Pyrrole~: deux signaux égaux
à 6.7 et \qty{6.2}{ppm} et un N--H large vers \qty{8.1}{ppm}. Furane~: deux signaux égaux à
7.4 et \qty{6.4}{ppm}, sans N--H. Thiophène~: deux signaux égaux à 7.33 et \qty{7.12}{ppm}, proches
l'un de l'autre, entre ceux du furane.
\end{solution}

\begin{solution}{pb:b3:heterocycles:1}
\textbf{1.} $\ce{C6H5NHNH2 + (CH3)2CO -> C6H5NHN=C(CH3)2 + H2O}$.

\textbf{2.} L'azote de l'amine s'additionne sur le carbone du carbonyle,
puis de l'eau est perdue, comme lors de la formation d'une imine~; l'acide
active le groupe carbonyle et transforme le OH de l'adduit en groupe
partant.

\textbf{3.} Le $\ce{NH2}$ terminal~: le doublet de l'autre azote est conjugué avec le
cycle phényle et plus encombré.

\textbf{4.} $\qty{108.0}{g/mol}$~; $\qty{10.0}{g}/\qty{108.0}{g/mol} = \qty{0.0926}{mol}$.

\textbf{5.} L'hydrazone $\ce{C9H12N2}$, \qty{148.0}{g/mol}~: $\qty{0.0926}{mol} \times \qty{148.0}{g/mol} = \qty{13.7}{g}$.

\textbf{6.} $\ce{C6H5-NH-NH-C(CH3)=CH2}$.

\textbf{7.} Le carbone \textit{ortho} du cycle, le carbone \textit{ipso}, les deux azotes, le carbone de
l'hydrazone et le $\ce{CH2}$ terminal.

\textbf{8.} La liaison N--N se rompt~; une liaison C--C se forme entre le carbone \textit{ortho} et le $\ce{CH2}$.

\textbf{9.} Six électrons, toutes les composantes \omterm{def:b3:pericyclic:faciality}{suprafaciales}~: thermiquement permise.

\textbf{10.} En tant qu'acide de Lewis, il catalyse la tautomérie, se lie à
un azote et affaiblit la liaison N--N, et aide l'ammoniac à partir.

\textbf{11.} Le carbone \textit{ortho} est devenu $sp^3$, avec un hydrogène~: la perte de ce proton
restaure le cycle benzénique, une forte force motrice.

\textbf{12.} Le $\ce{NH2}$ d'aniline formé attaque le carbone de l'imine, ce qui ferme le
cycle à cinq chaînons (une 2-aminoindoline).

\textbf{13.} L'ammoniac~; la formation du cycle pyrrole aromatique de l'indole entraîne l'étape.

\textbf{14.} $\ce{C9H12N2 -> C9H9N + NH3}$.

\textbf{15.} \qty{131.0}{g/mol}.

\textbf{16.} $\qty{0.0926}{mol} \times \qty{131.0}{g/mol} = \qty{12.1}{g}$.

\textbf{17.} En C3, la position la plus nucléophile de l'indole, libre
ici~; une attaque en ce point garde intact le cycle benzénique dans
l'intermédiaire.

\textbf{18.} $\ce{C6H5NHNH-C(=CH2)-CH2CH3}$ (terminale) et $\ce{C6H5NHNH-C(CH3)=CH-CH3}$ (interne).

\textbf{19.} 2-Éthylindole et 2,3-diméthylindole.

\textbf{20.} L'interne est la plus substituée~; les acides forts augmentent
la part de la terminale.

\textbf{21.} Le 2,3-diméthylindole, par l'ène-hydrazine la plus stable, la
plus substituée.

\textbf{22.} L'acétaldéhyde se condense sur lui-même et la voie par
l'ène-hydrazine donne de faibles rendements~; on prépare plutôt l'indole à
partir de l'hydrazone de l'acide pyruvique, qui donne l'acide
indole-2-carboxylique, ensuite décarboxylé.

\textbf{23.} Un mélange de régioisomères doit être séparé, ce qui coûte du
rendement~; on préfère une cétone symétrique ou une voie qui fixe la
régiochimie.

\textbf{24.} 10.0 g de phénylhydrazine peuvent donner au plus \qty{12.1}{g} de 2-méthylindole.
\end{solution}
